\chapter{Our Approach}

This chapter presents the methodological framework and empirical evaluation of the Retrieval-Augmented Generation (RAG) pipelines developed for this project. We investigate the efficacy of various retrieval strategies, ranging from lexical matching to advanced late-interaction neural architectures.

\section{Experiments on Embeddings and Retrieval Methods}

\subsection{Experimental Methodology}

\textbf{Dataset and Computing Environment}
To benchmark the performance of different retrieval architectures, we utilized the \textbf{Microsoft MS MARCO} (Machine Reading Comprehension) dataset\footnote{\url{https://huggingface.co/datasets/microsoft/ms_marco}}. This dataset is a standard benchmark for information retrieval tasks.

Due to computational constraints and to facilitate rapid iterative testing, we curated a representative subset of the data:
\begin{itemize}
    \item \textbf{Corpus Size:} 50,000 documents (reduced from the original $\approx$1 million).
    \item \textbf{Test Queries:} 556 distinct evaluation queries.
    \item \textbf{Hardware Infrastructure:} All experiments were conducted on Google Colab utilizing a single \textbf{Tesla T4 GPU}.
\end{itemize}

\textbf{Scoring Formulations}
Different pipelines utilize distinct mathematical formulations to calculate the relevance score $S(Q, D)$ between a query $Q$ and a document $D$.

\textbf{1. Lexical Scoring (BM25)}

For the sparse retrieval baseline, we employed the BM25 algorithm. This probabilistic model estimates relevance based on term frequency (TF) and inverse document frequency (IDF), normalizing for document length:
\begin{equation}
    \text{Score}_{\text{BM25}}(Q, D) = \sum_{t \in Q} \text{IDF}(t) \cdot \frac{\text{TF}(t, D) \cdot (k_1 + 1)}{\text{TF}(t, D) + k_1 \cdot \left(1 - b + b \cdot \frac{|D|}{\text{avgdl}}\right)}
\end{equation}
This method focuses on exact keyword matching and is highly effective for queries containing specific identifiers (e.g., proper nouns), though it lacks semantic understanding.

\textbf{2. Semantic Scoring (Dense Retrieval)}

For dense retrieval (e.g., MiniLM, BGE), we encode both query and document into fixed-size vectors ($\vec{V}_Q, \vec{V}_D$) and calculate the Cosine Similarity. This captures semantic meaning beyond exact keyword overlap:
\begin{equation}
    \text{Score}_{\text{Dense}}(Q, D) = \frac{\vec{V}_{Q} \cdot \vec{V}_{D}}{|\vec{V}_{Q}| \cdot |\vec{V}_{D}|}
\end{equation}
The score range is $[-1, 1]$, representing the semantic proximity in the vector space.

\textbf{3. Late Interaction Scoring (ColBERT/ColQwen)}

Unlike dense retrievers that compress a document into a single vector, the Late Interaction architecture (used in ColBERTv2) encodes queries and documents into bags of embedding vectors. The relevance score is the sum of the maximum similarity (MaxSim) for each query token against all document tokens:
\begin{equation}
    \text{Score}_{\text{MaxSim}}(Q, D) = \sum_{q \in \vec{Q}} \max_{d \in \vec{D}} (\vec{q} \cdot \vec{d})
\end{equation}
This allows for fine-grained matching, ensuring that specific details in a document are not lost during vector compression.

\textbf{Evaluation Metrics}
To quantitatively assess the pipelines, we adhere to two standard information retrieval metrics:

\begin{enumerate}
    \item \textbf{Recall@k:} Measures the percentage of relevant documents successfully retrieved within the top $k$ results.
    \[
    \text{Recall@k} = \frac{| \{ \text{Relevant Docs} \} \cap \{ \text{Retrieved Top-k} \} |}{| \{ \text{Total Relevant Docs} \} |}
    \]
    \item \textbf{nDCG@k (Normalized Discounted Cumulative Gain):} This metric evaluates the \textit{ranking quality}. It rewards algorithms that place relevant documents at the top of the list.
    \[
    \text{nDCG@k} = \frac{DCG@k}{IDCG@k} \quad \text{where} \quad DCG@k = \sum_{i=1}^{k} \frac{2^{rel_i} - 1}{\log_2(i+1)}
    \]
\end{enumerate}

\subsection{Pipeline Analysis and Results}

We conducted a comprehensive evaluation of Lexical, Dense, Hybrid, and Reranking pipelines. The aggregate results are presented in Table \ref{tab:results}.

\begin{table}[H]
\centering
\resizebox{\textwidth}{!}{%
\begin{tabular}{@{}llccrr@{}}
\toprule
\textbf{Pipeline Strategy} & \textbf{Model / Configuration} & \textbf{Latency} & \textbf{nDCG@3} & \textbf{Recall@3} \\ \midrule
\multicolumn{5}{l}{\textbf{Baseline Retrievers (No Rerank)}} \\
Sparse & BM25 (rank-bm25) & 3m & 0.2821 & 0.3085 \\
Dense & MiniLM-L6 (Raw) & $\sim$1s & 0.8658 & 0.9125 \\
Dense & \textbf{BGE-Small-en-v1.5} & $\sim$1s & \textbf{0.8855} & \textbf{0.9257} \\ \midrule
\multicolumn{5}{l}{\textbf{Hybrid Retrievers (Reciprocal Rank Fusion)}} \\
Hybrid & BM25 + MiniLM ($k=60$) & $\sim$1s & 0.5786 & 0.6885 \\
Hybrid & BM25 + BGE ($k=60$) & $\sim$1s & 0.5830 & 0.6787 \\ \midrule
\multicolumn{5}{l}{\textbf{Reranking Pipelines (Retrieve + Cross-Encoder)}} \\
Retrieve + Rerank & BM25 + BGE-Large & 1m 43s & 0.4469 & 0.4592 \\
Retrieve + Rerank & MiniLM + BGE-Large & 1m 43s & 0.9196 & 0.9526 \\
Retrieve + Rerank & \textbf{BGE-Small + BGE-Large} & 1m 43s & \textbf{0.9248} & \textbf{0.9598} \\ \midrule
\multicolumn{5}{l}{\textbf{Late Interaction}} \\
Late Interaction & \textbf{ColBERTv2 (RAGatouille)} & 6.5s & \textbf{0.9638} & \textbf{0.9412} \\ \bottomrule
\end{tabular}%
}
\caption{Comparative performance of Retrieval and Reranking pipelines on the MS MARCO subset (50k docs, 556 queries).}
\label{tab:results}
\end{table}

\textbf{Question 1: Is Hybrid Retrieval superior to single-mode retrieval?}
Our experiments demonstrate that Hybrid Retrieval (combining BM25 and Dense vectors via Reciprocal Rank Fusion) offers a significant improvement over purely lexical methods, providing a balanced approach to robustness.

\begin{itemize}
    \item \textbf{Limitations of BM25:} The raw BM25 baseline yielded a low nDCG@3 of $0.2821$. This confirms that the MS MARCO dataset relies heavily on semantic understanding rather than exact keyword matches, causing BM25 to fail when vocabulary mismatches occur.
    \item \textbf{The Hybrid Improvement:} By fusing BM25 with Dense retrieval (Hybrid BGE), the nDCG@3 score nearly doubled to $0.5830$.
    \item \textbf{Significance:} While pure dense models performed exceptionally well on this specific dataset ($0.8855$), relying solely on dense retrieval in production carries the risk of "semantic hallucination" (retrieving conceptually similar but factually incorrect documents). Hybrid retrieval mitigates this by anchoring the semantic search with the keyword precision of BM25. It provides a "safety net" that ensures documents containing critical entity names are not discarded, making it a more robust strategy for real-world applications where query intent varies.
\end{itemize}

\textbf{Question 2: Does the addition of a Reranker significantly improve ranking quality?}
The results strongly support the hypothesis that a two-stage "Retrieve-then-Rerank" pipeline outperforms single-stage retrieval.

\begin{itemize}
    \item \textbf{Architectural Advantage:} We employed Cross-Encoders (specifically \textbf{BAAI/bge-reranker-large}) as the second stage. Unlike bi-encoders, cross-encoders process the query and document simultaneously in the Transformer, capturing deep interaction signals.
    \item \textbf{Performance Gain:} Applying the BGE-Large reranker to the BGE-Small retriever candidates improved nDCG@3 from $0.8855$ to \textbf{0.9248} and Recall@3 from $0.9257$ to \textbf{0.9598}.
    \item \textbf{Model Comparison:} We observed that larger rerankers yield better quality. The \textbf{bge-reranker-large} consistently outperformed the smaller \textbf{ms-marco-MiniLM-L-12} (e.g., 0.9248 vs 0.9107). However, this comes with a trade-off: the large reranker is approximately $6\times$ slower ($1m43s$ vs $17s$).
    \item \textbf{Conclusion:} For applications prioritizing accuracy, the combination of a high-performance dense retriever (BGE-Small) and a large cross-encoder (BGE-Large) is the optimal two-stage configuration.
\end{itemize}

\textbf{Question 3: Is ColBERT (Late Interaction) the most effective approach?}
Our data indicates that ColBERTv2 provides the highest retrieval quality, surpassing even the best Retrieve-then-Rerank pipelines.

\begin{itemize}
    \item \textbf{Superior Metrics:} ColBERTv2 achieved the highest \textbf{nDCG@3 of 0.9638}, outperforming the best reranking pipeline ($0.9248$).
    \item \textbf{Solving the Information Bottleneck:} Standard two-stage pipelines suffer from a bottleneck: if the first-stage retriever (e.g., BGE-Small) fails to find the relevant document in the top 30 candidates, the Reranker never sees it. ColBERT avoids this by using a "Late Interaction" mechanism. It calculates the MaxSim between query tokens and document tokens directly.
    \item \textbf{Granularity:} As shown in Equation (3.3), ColBERT matches at the token level. This allows it to identify a relevant document even if only a small paragraph within it matches the query—a nuance often lost when compressing a whole document into a single vector (Dense Retrieval).
    \item \textbf{Operational Note:} It is worth noting that while ColBERT's retrieval time is efficient ($6.5s$ on GPU), the indexing process is resource-intensive. In our environment, indexing required $\approx35$ minutes on CPU\footnote{Due to FAISS-GPU limitations on the Colab instance, indexing fell back to CPU.} compared to mere seconds/minutes for standard dense models. Despite this setup cost, ColBERT represents the state-of-the-art for retrieval accuracy on this dataset.
\end{itemize}

\section{Justification for Retrieval-Augmented Generation (RAG)}

While Large Language Models (LLMs) possess remarkable generative capabilities, they suffer from inherent limitations such as knowledge cutoffs, hallucinations, and a lack of access to proprietary data. To address these issues, several methodologies exist. This section analyzes alternative approaches - Fine-Tuning, Semantic Search, Prompt Engineering, and Domain-Specific Pre-training - and justifies the selection of RAG as the superior architectural choice for this project.

\subsection{Comparative Analysis of Alternatives}

\textbf{RAG vs. Fine-Tuning}
Fine-tuning involves retraining a pre-trained LLM on a specific dataset to adjust its internal parameters (weights). While this method tailors the model to a specific domain style, it was rejected for the following reasons:

\begin{itemize}
    \item \textbf{Static Knowledge Base:} Fine-tuned models immediately become outdated upon training completion. Updating the knowledge requires a computationally expensive retraining process. In contrast, RAG allows for \textit{dynamic updates} - simply adding a document to the database makes it immediately retrievable without modifying the model weights.
    \item \textbf{Hallucination Risk:} Fine-tuning increases the probability of the model adopting a specific tone, but it does not guarantee factual accuracy. The model can still hallucinate information effectively. RAG mitigates this by grounding the generation in retrieved evidence.
    \item \textbf{Explainability:} A fine-tuned model outputs an answer based on "black box" internal weights. RAG provides citation and source attribution (retrieved chunks), which is critical for verifying accuracy in our application.
\end{itemize}

\textbf{RAG vs. Prompt Engineering (Context Stuffing)}
Prompt engineering involves optimizing the textual input to guide the LLM's output, often by manually pasting context. While simple, it is insufficient for our scale:

\begin{itemize}
    \item \textbf{Context Window Limitations:} Although context windows are growing (e.g., 128k tokens), it is infeasible to pass an entire corpus (e.g., 50,000 documents) into a single prompt. RAG acts as a filter, selecting only the most relevant snippets to fit within the model's constraints.
    \item \textbf{Performance Degradation:} Studies show that as the context length increases, the model's ability to reason over the data ("Lost in the Middle" phenomenon) decreases. RAG ensures the model focuses only on highly relevant, condensed information.
\end{itemize}

\textbf{RAG vs. Semantic Search (Traditional Information Retrieval)}
Semantic search utilizes vector databases to retrieve documents based on meaning rather than keywords. However, it functions solely as a retrieval engine, not a generation engine.

\begin{itemize}
    \item \textbf{Synthesis vs. Retrieval:} Semantic search returns a list of document links or snippets. The user is then forced to read and synthesize the answer manually. RAG automates this cognitive load by employing the LLM to read the retrieved documents and generate a coherent, natural language answer.
    \item \textbf{Completeness:} A user's query might require synthesizing information from three different documents. Semantic search provides the three documents; RAG provides the unified answer derived from them.
\end{itemize}

\textbf{RAG vs. Domain-Specific Pre-training}
Pre-training involves training a model from scratch or continuing pre-training on a massive domain-specific corpus.

\begin{itemize}
    \item \textbf{Computational Cost:} This approach is prohibitively expensive, requiring massive GPU clusters and weeks of training time.
    \item \textbf{Flexibility:} Like fine-tuning, a pre-trained model fails in dynamic environments where data changes daily. RAG decouples the \textit{reasoning engine} (the LLM) from the \textit{knowledge base} (the Vector DB), allowing us to swap or update the knowledge base instantly without touching the neural network.
\end{itemize}

\subsection{Summary of Trade-offs}

The following table summarizes the key architectural decisions that led to the adoption of RAG.

\begin{table}[h]
\centering
\resizebox{\textwidth}{!}{%
\begin{tabular}{@{}lcccc@{}}
\toprule
\textbf{Feature} & \textbf{RAG (Ours)} & \textbf{Fine-Tuning} & \textbf{Prompt Engineering} & \textbf{Semantic Search} \\ \midrule
\textbf{Knowledge Update} & \textit{Real-time (Dynamic)} & Slow (Retraining) & Manual & Real-time \\
\textbf{Hallucination Risk} & Low (Grounded) & High & Medium & N/A (No Generation) \\
\textbf{Context Capacity} & Infinite (via Retrieval) & Limited by Training & Limited by Window & Infinite \\
\textbf{Computational Cost} & Low (Inference only) & High (Training) & Very Low & Low \\
\textbf{Answer Synthesis} & \textbf{Yes} & Yes & Yes & \textbf{No} \\ \bottomrule
\end{tabular}%
}
\caption{Comparison of RAG against alternative LLM enhancement strategies.}
\label{tab:rag_comparison}
\end{table}

\textbf{Conclusion:} 
RAG is chosen because it offers the optimal balance of accuracy, cost-efficiency, and dynamic adaptability. It leverages the reasoning power of pre-trained LLMs while overcoming their memory limitations through external retrieval, making it the most suitable architecture for a document-based Question Answering system.

\newpage